\documentclass{article}
\usepackage{amsmath}
\usepackage{amssymb}
\usepackage{amsthm}
\usepackage{mathtools}
\usepackage{geometry}
\usepackage{hyperref}

\newcommand{\rk}{\operatorname{rank}}
\newcommand{\GL}{\operatorname{GL}}
\newcommand{\R}{\mathbb{R}}
\newcommand{\Mat}{\mathrm{Mat}}

\theoremstyle{definition}
\newtheorem{theorem}{Theorem}[section]
\newtheorem{proposition}[theorem]{Proposition}
\newtheorem{lemma}[theorem]{Lemma}
\newtheorem{corollary}[theorem]{Corollary}
\newtheorem{definition}[theorem]{Definition}
\newtheorem{remark}[theorem]{Remark}
\newtheorem{conjecture}[theorem]{Conjecture}

\title{Task 5: Statement of the Second Conjecture}
\date{}

\begin{document}
\maketitle

\section{Goal}

This file states the second conjecture in its final form.

The conjecture should be read as the module-theoretic companion to the first
conjecture:

\begin{itemize}
\item the first conjecture identifies the correct orbit relation between the
  weight-space cyclic locus and the critical locus;
\item the second conjecture identifies the combinatorial labels that are carried
  by those orbits.
\end{itemize}

\section{The module attached to a point of the cyclic locus}

\begin{definition}[The $A_L^{\mathrm{wt}}$-module attached to $W$]
Fix the closing map
\[
W_0=\Sigma_{XY}P_Q.
\]
For
\[
W=(W_1,\dots,W_L)\in \Gamma_{d,S},
\]
define the associated quiver representation
\[
\phi(W):=(W_0,W_1,\dots,W_L).
\]
By Tasks 0 and 1, this defines an $A_L^{\mathrm{wt}}$-module, which we denote by
\[
M(W).
\]
\end{definition}

\begin{definition}[Multiplicity pattern of a cyclic point]
Using the interval decomposition from Task 4, write
\[
M(W)
\cong
\bigoplus_{(i,j)\in \mathcal I_L^{\mathrm{wt}}}
M(i,j)^{\oplus m_W(i,j)}.
\]
The collection
\[
m_W:=(m_W(i,j))_{(i,j)\in \mathcal I_L^{\mathrm{wt}}}
\]
is called the multiplicity pattern of $W$.
\end{definition}

\section{Second conjecture}

\begin{conjecture}[Second conjecture: admissible-interval classification in weight space]
Let
\[
A_L^{\mathrm{wt}}=kC_{L+1}/I_L^{\mathrm{wt}}
\]
be the weight-space cyclic algebra with
\[
I_L^{\mathrm{wt}}=\langle p_1,\dots,p_L\rangle.
\]
Then the combinatorics of the direct weight-space cyclic locus is controlled by
the admissible interval modules
\[
\{M(i,j)\}_{(i,j)\in \mathcal I_L^{\mathrm{wt}}},
\qquad
\mathcal I_L^{\mathrm{wt}}
=
\{(i,j)\mid 0\le i\le j\le L\}
\cup
\{(i,j)\mid 1\le i\le L,\ 0\le j\le i-2\}.
\]

More precisely, the following two statements are equivalent formulations of the
same conjectural picture.

\medskip
\noindent
\textbf{Abstract-module form.}
Every finite-dimensional $A_L^{\mathrm{wt}}$-module of dimension vector
\[
d=(d_0,\dots,d_L)
\]
admits a unique decomposition
\[
M
\cong
\bigoplus_{0\le i\le j\le L} M(i,j)^{\oplus m(i,j)}
\oplus
\bigoplus_{\substack{1\le i\le L\\0\le j\le i-2}}
M(i,j)^{\oplus m(i,j)},
\]
with multiplicities
\[
(m(i,j))\in \mathcal M_d^+(A_L^{\mathrm{wt}}).
\]

\medskip
\noindent
\textbf{Weight-space form.}
Every point
\[
W\in \Gamma_{d,S}
\]
with fixed closing map $W_0=\Sigma_{XY}P_Q$ determines a unique multiplicity
pattern
\[
m_W\in \mathcal M_d^+(A_L^{\mathrm{wt}})
\]
through the decomposition of the associated module $M(W)$.
These multiplicity patterns are the combinatorial labels of the cyclic locus and
should be the labels transferred to the critical locus through the orbit
correspondence of the first conjecture.
\end{conjecture}

\section{Expected stratification language}

\begin{definition}[Cyclic multiplicity stratum]
For
\[
m\in \mathcal M_d^+(A_L^{\mathrm{wt}}),
\]
define
\[
\Gamma_{d,S}[m]
:=
\left\{
W\in \Gamma_{d,S}
\;\middle|\;
M(W)\cong
\bigoplus_{(i,j)\in \mathcal I_L^{\mathrm{wt}}}
M(i,j)^{\oplus m(i,j)}
\right\}.
\]
\end{definition}

\begin{remark}
The conjectural role of the second conjecture is that the sets
\[
\Gamma_{d,S}[m]
\]
give the correct combinatorial partition of the cyclic locus. Once the
first-conjecture orbit correspondence is available, the same labels should
induce the corresponding partition of the critical locus.
\end{remark}

\begin{remark}[What is deliberately excluded here]
This conjecture does \emph{not} yet include:

\begin{itemize}
\item Ext formulas;
\item codimension formulas;
\item orbit-closure relations;
\item rank-selection rules inside the fixed-closing-arrow slice.
\end{itemize}

Those belong to later parts of the project.
\end{remark}

\section{Deliverables}

This task is complete once the file contains:

\begin{itemize}
\item the module $M(W)$ attached to a point of $\Gamma_{d,S}$;
\item the exact admissible summation ranges for the decomposition;
\item the second conjecture in abstract-module form;
\item the same conjecture in weight-space form;
\item the definition of the multiplicity strata $\Gamma_{d,S}[m]$.
\end{itemize}

\end{document}
